\section{Discussion: A Scientific Identification Gate}
\label{sec:gate}

The audit in \cref{sec:results} is not a catalogue of incidental shortcomings; the
three findings share a common root. A pooled, linear, observational model of ambient
telemetry cannot, in principle, answer the interventional and apportionment questions
the original pipeline posed. We formalize this as a four-part \emph{identification
gate}: criteria that an observational air-quality analysis must pass before causal,
apportionment or policy language is warranted.

\subsection{Gate 1: structural absence of VOCs}
Ozone sensitivity is governed by the VOC-to-NO$_x$ ratio
\citep{sillman1999,monks2015ozone,seinfeld2016atmospheric}. Without continuous
hydrocarbon speciation---or an observational proxy such as the formaldehyde-to-NO$_2$
column ratio (HCHO/NO\textsubscript{2}) from satellite retrievals---the VOC-limited and
NO$_x$-limited regimes are empirically indistinguishable in these data. A model that
never observes VOCs cannot resolve which regime produced a given ozone value, and so
cannot attribute an ozone change to a precursor intervention. Our single recovered
edge, $\mathrm{NO}\rightarrow\mathrm{O_3}<0$, is consistent with this limitation: it is
the \emph{titration} term, the one piece of ozone chemistry that does not require
observing VOCs, and it must not be generalized to regional photochemical production.

\subsection{Gate 2: unobserved planetary-boundary-layer confounding}
The diurnal cycle of the planetary boundary layer (PBL) is a powerful \emph{unobserved
common cause}. At night and in winter the PBL collapses, trapping fresh emissions in a
shallow volume and simultaneously raising the measured concentrations of all primary,
co-emitted species; by day it deepens and dilutes them together. This single latent
variable induces strong positive correlations among co-emitted pollutants that have
nothing to do with any direct inter-pollutant causal mechanism, violating the causal
sufficiency assumption (A4) of \eqref{eq:sem}. The empirical signature is exactly what
we observe: the particulate and combustion edges of \cref{fig:dag} are stable in the
pollutant-only graph but collapse once meteorological proxies of mixing are included.
A spuriously ``discovered'' $\mathrm{CO}\rightarrow\mathrm{PM_{2.5}}$ edge is better read
as two species responding to a common, unmeasured mixing volume. Mitigation requires
either direct PBL-height observations or explicit meteorological deconfounding
\citep{grange2019metnorm,grange2018rf,chernozhukov2018double}.

\subsection{Gate 3: the concentration--emission fallacy}
This is the decisive gate for policy. The optimizer of \cref{lem:pareto} acts on
\emph{standardized ambient concentrations}: a recommendation such as
``$z_{\mathrm{NO_2}}=z_{\min}=-2\sigma$'' is a point on a correlational surface, not an
emission control. Two failures follow. First, there is no validated map from this
concentration vector back to a feasible set of sectoral emission reductions; the
inverse problem is precisely the ill-posed source attribution that
\eqref{eq:adr} would have to solve. Second, the search can select physically
impossible states: for a right-skewed, non-negative pollutant whose standardized mean
sits near zero, a target of $-2\sigma$ can correspond to a \emph{negative physical
mass concentration}. Formally, if $C=\mu+\sigma z$ with $\mu\lesssim 2\sigma$ (as holds
for several species in \cref{tab:summary}), then $z=-2$ implies $C<0$. An
``optimal policy'' that prescribes negative pollutant mass is a reductio: it exposes
that the objective never encoded emissions, feasibility, or benefit. Optimizing a
concentration response surface and reporting the result as a control strategy is a
category error, independent of which optimizer is used.

\subsection{Gate 4: leakage-aware evaluation}
Finally, no skill claim survives without an autocorrelation-aware protocol. As
\cref{tab:bench} shows, the same model and data yield $R^2$ from $0.41$ to $0.22$
depending solely on the split. Random $k$-fold over rows of spatio-temporal telemetry
is not a valid estimate of out-of-sample skill \citep{roberts2017cv,kaufman2012leakage};
blocked-temporal and leave-groups-out protocols, benchmarked against persistence and
climatology, are the minimum standard.

\subsection{Recommendations}
For researchers and regulators deploying observational environmental AI we recommend:
(i) report blocked and grouped validation with persistence and climatology baselines,
never a single random-split number; (ii) treat discovered graphs as hypotheses and test
every edge against meteorological confounding and bootstrap stability before any causal
language; (iii) verify measurement units and provenance against authoritative sensor
metadata, not value heuristics; (iv) reserve ``source apportionment'' for receptor
models \eqref{eq:pmf} with speciation and uncertainties; and (v) never present an
optimizer over ambient concentrations as an emission-control recommendation without a
validated, feasibility-constrained emissions-to-concentrations model.
